% !PRINTED 443
\weberpage{443}
% running head: § 122.  Irreducibility.  443

%<<CONT
of the field $\Omega$ (Bd.\,II, § 149). If $t$ denotes a variable, then $\Phi(t)$ is a divisor of the class function $H_m(t)$, and consequently the discriminant of $\Phi(t)$ is also a divisor of the discriminant of $H_m(t)$.

If $\xi$ is an arbitrary integer of the field $\mathfrak{K}$, then the relative trace with respect to $\Omega$ (Bd.\,II, § 175):
\[
\mathfrak{S}\left(\frac{\xi}{t-(k)}\,\Phi(t)\right) = \varphi(t)
\]
is an integral function of $t$, of degree at most $\nu-1$, whose coefficients are \gesperrt{integers} of $\Omega$. If we therefore let $t$ pass over into $(k)$, it follows that:
\weq{3}{\xi = \frac{\varphi(k)}{\Phi'(k)},}
and from this:
\weq{4}{b\xi = \beta_0 + \beta_1(k) + \beta_2(k)^2 + \cdots + \beta_{\nu-1}(k)^{\nu-1},}
in which $\beta_0, \beta_1, \ldots, \beta_{\nu-1}$ are \gesperrt{integers} in $\Omega$, and $b$ denotes an integral rational number, for which the discriminant of equation (2):
\[
b = N[\Phi'(k)]
\]
may be taken. The field discriminant then divides $b$, and \gesperrt{none of the primes not dividing $b$ is divisible in $\Omega$ or in $\mathfrak{K}$ by the square of a prime ideal} (Bd.\,II, § 174). The number $b$ does not depend on the particular number $\xi$.

An integer $\xi$ can be represented in the form (4) in only one way, since otherwise $(k)$ would satisfy an equation of lower than the $\nu$th degree, which contradicts the assumption; \gesperrt{and from this it follows that $\xi$ is divisible by an integer $\alpha$ of the field $\Omega$ only if all the coefficients $\beta_0, \beta_1, \ldots, \beta_{\nu-1}$ are divisible by $\alpha$.} For if one represents $\xi/\alpha$ in the form (4), the coefficients in this representation must likewise be integers.

If we replace $(k)$ in (3) by another root of equation (2), then $\xi$ passes over into a conjugate value, which is likewise an integer.

2. By \gesperrt{Fermat's} theorem we have, for every arbitrary prime $p$, when $a$ denotes an integral rational number:
\weq{5}{a^p \equiv a \pmod{p}.}
